\section{Related Work}
\label{sec:related}

\subsection{Persona Prompting Effects on AI Output}

The use of persona descriptions to condition large language model behavior has a substantial empirical literature. Zhang et al.~\cite{zhang2018personas} demonstrate that injecting persona descriptions into dialogue systems improves response coherence and user engagement, establishing that LLMs benefit from explicit identity framing rather than generic instruction. Park et al.~\cite{park2023generative} extend this to generative agents with persistent identity representations, showing that memory-augmented personas can sustain coherent behavior across multi-turn social interactions. Dang et al.~\cite{dang2023personas} apply persona conditioning to creative writing assistance, finding that distinct persona instantiations produce meaningfully different stylistic outputs from the same underlying model, with implications for diversity in co-creative systems.

A consistent limitation across this prior literature is that persona is treated as a holistic unit. A role label is applied; the model responds from that role; output quality is attributed to the persona as a whole. The implicit theoretical assumption is that any adequate persona specification will activate the relevant knowledge and values the model associates with that identity, and that the primary design variable is \emph{which} persona to select rather than \emph{how} to specify it. Dang et al. acknowledge that different creative personas produce different outputs but do not ask which internal component of the persona description is responsible for particular output properties~\cite{dang2023personas}. Park et al. focus on persona \emph{persistence} over time rather than on persona \emph{composition}---the degree to which stable behavior across turns is attributable to specific profile elements versus the identity label as a gestalt~\cite{park2023generative}.

More recent work on AI reviewer personas has shown that the \emph{richness} of a persona specification, as distinct from its mere presence, materially affects output quality. \citet{dellaLibera2025panel} demonstrate that domain-specific profiles encoding named evaluative frameworks, characteristic concerns, and analytical vocabulary produce qualitatively different review feedback from generic expert labels---not primarily in score magnitude but in the specificity, transferability, and actionability of the frameworks introduced. Crucially, however, that work operates on the \emph{reviewer} side of creative evaluation: it asks how persona depth affects the quality of assessment, not how persona identity affects the quality of creative generation. The present paper fills the complementary gap: we ask which components of a creator persona are responsible for creative output quality, using the same domain (puzzle hunt design) and the same evaluation infrastructure.

To our knowledge, no prior work has decomposed a persona prompt into constituent cognitive traits, isolated those traits experimentally, and identified which subset causally drives observed quality outcomes. The present paper introduces this paradigm.

\subsection{Creative Cognition and Expertise}

The relationship between domain expertise and creative output is contested in creativity research, and the artist-versus-puzzle-designer finding lands directly in the middle of that debate.

Csikszentmihalyi's systems model~\cite{csikszentmihalyi1996creativity} situates creative contribution at the intersection of three entities: the individual (person), the symbolic domain (accumulated knowledge and conventions), and the social field (gatekeepers and evaluators who recognize genuine novelty). On this account, domain expertise is necessary but not sufficient for creativity---the individual must internalize the domain's conventions deeply enough to transform them, rather than merely reproduce them. The artist persona finding is consistent with this model: the puzzle designer who produces an acrostic extraction table has internalized the domain's conventions but has not transformed them. The artist, lacking those conventions, designs from first principles and produces a structurally novel mechanism.

Weisberg's work on creative expertise~\cite{weisberg2006creativity} complicates this picture. Weisberg argues against the romantic view of creativity as transcendent novelty, proposing instead that creative advances are produced by deep domain knowledge combined with incremental variation. On this account, the artist's advantage should be impossible: without puzzle domain knowledge, the artist cannot produce meaningful novelty. That the artist nonetheless produces the study's only novel mechanism suggests that Weisberg's account, while broadly accurate for expert practitioners, underweights the contribution of \emph{design stance}---the cognitive orientation from which the creator approaches the task---relative to domain content. The trait decomposition reported here provides a mechanistic account of what ``design stance'' consists of and which of its components are load-bearing.

Sawyer's work on improvisation and creative constraint~\cite{sawyer2003improvisation} is relevant to understanding why artistic orientation may outperform domain expertise for generative tasks. Sawyer argues that experienced improvisers work within structures they have internalized so deeply that the structures become invisible, freeing cognitive resources for experience-oriented decisions (``what should happen next for the audience?'') rather than rule-compliance decisions (``what is allowed next given the grammar?''). The artist persona activates this same reorientation: the trait-isolation results show that experience-first thinking (T1) is the primary driver of quality, while systematic thinking (T6)---the trait most analogous to rule-grammar compliance---produces the lowest isolation score in the study (22.3/45, below the puzzle designer baseline of 24.7/45). Cognitive freedom from convention, in this account, is not ignorance but a different allocation of attention.

\subsection{Experience-Design and Human-Centered Creation}

The experience-first trait (T1)---``you always think first about how the other person will experience what you make''---is the operationalization of a principle with a rich theoretical basis in human-centered design. Norman's foundational work~\cite{norman2013design} establishes that the user's experience of an artifact is not a downstream property of good engineering but its primary constraint: design that begins with structure and adds usability as an afterthought produces categorically inferior results compared to design that begins with the user's mental model and builds structure to serve it. Hassenzahl's theory of experience design~\cite{hassenzahl2010experience} extends this principle from usability to hedonic quality, arguing that the most important design question is not ``what does this product do?'' but ``what does it feel like to use it?''---a reframing that shifts the designer's primary object of attention from feature specification to experiential arc.

In the puzzle design literature, this principle is named explicitly by Thomas Snyder, whose review frameworks appear throughout the C8 panel evaluation data. Snyder's central construction criterion---``would you want to solve this?''---is experience-first design applied to puzzle authoring~\cite{snyder2019craftsmanship}. A constructor who asks ``what mechanism yields TOWER?'' designs differently from one who asks ``what would it feel like to discover that the answer is TOWER?'' The damaged-inscription mechanism in AA is the product of the second question: the mechanism \emph{is} the experience of recognition---the solver sees a gap in a word, reconstructs the missing letter from domain knowledge, and experiences a moment of identification that is simultaneously archaeological and self-confirming. This is not a mechanism overlaid on an experience; it is a mechanism that \emph{is} the experience.

The leave-one-out results confirm that experience-first thinking is not merely one quality-amplifying factor among six---it is the architectural foundation of the puzzle as a solvable artifact. Without T1, the resulting ``puzzle'' (L-T1, scored 12.3/45) is a list of AoE2 terms with no clue layer, no mechanism, and no solver task. The model without the experience-first orientation does not produce a worse puzzle; it produces an object that is not a puzzle at all.

\subsection{Minimalism and Reduction in Design Practice}

The elegance drive trait (T4)---``you remove anything that isn't necessary; you value spare, clean form above all''---is the operationalization of a design principle with as strong a theoretical pedigree as experience-first thinking. Dieter Rams articulated this principle as the tenth of his ten principles for good design: ``Good design is as little design as possible''~\cite{rams1976design}. In Rams's framework, reduction is not aesthetic preference but epistemic discipline: excess material obscures the relationship between form and function, making it harder for the user to understand what the artifact does and how to use it. Antoine de Saint-Exup\'ery's often-cited formulation---``Perfection is achieved not when there is nothing more to add, but when there is nothing left to take away''---makes the same point as a phenomenological principle: the experience of a well-designed artifact is the experience of inevitability, and inevitability is impossible in the presence of excess.

The leave-one-out result for elegance drive (L-T4, drop of $-$11.7 points, below-threshold at 30.0/45) provides empirical support for these design principles in a novel experimental context. The degraded condition does not merely score lower---it self-sabotages in a specific and diagnostic way. The closing note in L-T4 telegraphs the answer TOWER before the solver has had the opportunity to find it, converting the extraction from a discovery into a confirmation of what the solver was already told to expect. The unedited puzzle undermines its own mechanism. This is precisely the failure mode Rams predicts: excess material does not merely add noise; it corrupts the signal relationship between the designer's intention and the solver's experience.

Experimental aesthetics literature provides complementary evidence. Reber et al.~\cite{reber2004processing} demonstrate that perceptual fluency---the ease with which a stimulus can be processed---produces aesthetic pleasure, and that sparse form generally produces higher fluency than cluttered form. The elegant puzzle is not merely more pleasing; it is more legible, and that legibility is itself a quality of the solver's experience. The elegance drive trait operationalizes this finding as an authoring constraint: before calling work done, remove anything that reduces the clarity of the signal relationship between clue and answer.

\subsection{Prior Work on Minimal Sufficient Conditions in Creative Systems}

The Phase~3 reconstruction experiment---starting from a null baseline and adding traits until the quality threshold is crossed---is an application of the \emph{minimum sufficient set} paradigm to creative generation prompting. This paradigm has precedents in adjacent literatures. In automated theorem proving, minimal proof certificate identification has been studied as a problem of identifying which axioms and inference steps are load-bearing and which are redundant~\cite{harrison2009minproof}. In feature selection for machine learning, ablation-based subset identification is standard methodology. In human-computer interaction, the minimum viable prototype paradigm applies a related logic to product development~\cite{ries2011lean}.

What is novel in the present application is the identification of a minimum sufficient cognitive trait set for an open-ended creative generation task, using a human expert panel as the quality oracle. This differs from minimum viable prototype work in that the quality criterion is evaluative (expert panel score) rather than functional (does the feature work), and from machine learning ablation in that the ``features'' being tested are cognitive orientations rather than measurable input variables. The result---that T1 + T4 + T5 constitute the minimum sufficient set for passing-threshold puzzle authoring quality, with each Tier~3 trait contributing 4--8 additional quality points---is, to our knowledge, the first empirically identified minimum sufficient trait set for a persona prompting effect in any creative domain.

The gap this work fills is therefore threefold: no prior study decomposes professional identity prompts into constituent cognitive traits, no prior study identifies the minimum sufficient trait subset empirically, and no prior study demonstrates that domain-naive artistic orientation outperforms domain-specific expertise labels for creative generation tasks.
